\section{Entidades y atributos}

Esta sección registra las entidades y los atributos que se desprenden del caso
de uso oficial de PuellaGame. Los nombres de los identificadores se mantienen
consistentes con el resto del modelo: cada identificador permite distinguir de
manera única a una entidad y no se confunde con un dato descriptivo.

\subsection{Convenciones para los atributos}

Los dominios que se presentan son dominios conceptuales; no representan aún
tipos específicos de un sistema gestor de bases de datos. Cuando el caso de
uso lo indica, se conserva la naturaleza del atributo:

\begin{itemize}[leftmargin=*]
  \item un atributo compuesto se divide en componentes con significado
        propio, como el nombre completo y la dirección;
  \item un atributo multivalorado puede tener varios valores para una misma
        entidad, como los teléfonos, los correos y los horarios;
  \item un atributo derivado se obtiene a partir de otros datos, como la edad
        a partir de la fecha de nacimiento, la categoría VIP a partir de las
        visitas y la bonificación VIP a partir de la categoría vigente del
        cliente.
\end{itemize}

Las cantidades monetarias se expresan en pesos mexicanos (MXN), las fechas y
horas corresponden al momento registrado por la operación, y los dominios
enumerados conservan literalmente las opciones indicadas en el caso de uso.

\subsection{Entidades del servicio}

\begin{center}
\footnotesize
\renewcommand{\arraystretch}{0.8}
\begin{tabularx}{\textwidth}{@{}>{\bfseries\color{unamblue}}p{0.19\textwidth}p{0.28\textwidth}X@{}}
\toprule
Entidad & Atributo & Dominio y caracterización \\
\midrule
Sucursal & \texttt{idSucursal} & Identificador único; entero positivo. \\
 & \texttt{nombre} & Cadena de texto no vacía. \\
 & \texttt{dirección} & Atributo compuesto por \texttt{calle},
   \texttt{númeroInterior}, \texttt{númeroExterior}, \texttt{colonia} y
   \texttt{estado}; cada componente conserva el dominio propio de una
   dirección. \\
 & \texttt{teléfono} & Cadena con formato de número telefónico. \\
 & \texttt{horario} & Atributo multivalorado; conjunto de intervalos de atención
   expresados mediante hora inicial y hora final. \\
\midrule
Cliente & \texttt{idCliente} & Identificador único; entero positivo. \\
 & \texttt{nombreCompleto} & Atributo compuesto por \texttt{nombre},
   \texttt{apellidoPaterno} y \texttt{apellidoMaterno}; cadenas de texto. \\
 & \texttt{fechaNacimiento} & Fecha válida de nacimiento. \\
 & \texttt{edad} & Atributo derivado; número entero no negativo calculado a
   partir de \texttt{fechaNacimiento}. \\
 & \texttt{sexo} & Enumeración: \texttt{masculino}, \texttt{femenino} o
   \texttt{no binario}. \\
 & \texttt{correoElectrónico} & Atributo multivalorado; conjunto de cadenas
   con formato de correo electrónico. \\
 & \texttt{teléfono} & Atributo multivalorado; conjunto de cadenas con formato
   de número telefónico. \\
 & \texttt{categoríaVIP} & Atributo derivado; indica si el cliente tiene la
   categoría VIP conforme al número de visitas registrado. \\
\midrule
Tarjeta & \texttt{numeroTarjeta} & Identificador único; cadena o número de
   tarjeta no repetido. \\
 & \texttt{costoAdquisición} & Cantidad monetaria en MXN; el caso de uso fija
   el costo de adquisición en \$20.00. \\
 & \texttt{fechaCompra} & Fecha en la que se adquirió la tarjeta. \\
 & \texttt{saldo} & Cantidad monetaria en MXN; valor no negativo. \\
 & \texttt{puntosAcumulados} & Número entero no negativo. \\
 & \texttt{estado} & Enumeración que distingue, al menos, los estados
   \texttt{activa} e \texttt{inactiva}; permite expresar la regla de una sola
   tarjeta activa por cliente. \\
\midrule
Visita & \texttt{idVisita} & Identificador único; entero positivo. \\
 & \texttt{fecha} & Fecha en la que el jugador utilizó su tarjeta en una
   sucursal. \\
\midrule
Recarga & \texttt{idRecarga} & Identificador único; entero positivo. \\
 & \texttt{monto} & Cantidad monetaria en MXN, positiva y múltiplo de diez. \\
 & \texttt{fecha} & Fecha en la que se realizó la recarga. \\
 & \texttt{hora} & Hora en la que se realizó la recarga. \\
 & \texttt{métodoAbono} & Enumeración: \texttt{efectivo}, \texttt{tarjeta de
   crédito o débito} o \texttt{PayPal}. \\
\midrule
Canje & \texttt{idCanje} & Identificador único; entero positivo. \\
 & \texttt{fecha} & Fecha en la que se realizó el canje. \\
 & \texttt{hora} & Hora en la que se realizó el canje. \\
 & \texttt{cantidadPremios} & Número entero positivo. \\
 & \texttt{puntosUtilizados} & Número entero de puntos, con mínimo de veinte. \\
\bottomrule
\end{tabularx}
\end{center}

\Needspace{5\baselineskip}
El costo de adquisición y la fecha de compra de la tarjeta se representan
explícitamente porque aparecen como atributos en el diagrama. La recarga
inicial de \$10.00 MXN sigue siendo una regla del negocio de la operación de
venta y recarga. La categoría VIP de \texttt{Cliente} y la bonificación VIP de
\texttt{Partida} se marcan como derivadas; su justificación y sus restricciones
se documentan en las secciones correspondientes.

\subsection{Entidades de juegos y operación}

\begin{center}
\small
\begin{tabularx}{\textwidth}{@{}>{\bfseries\color{unamblue}}p{0.19\textwidth}p{0.28\textwidth}X@{}}
\toprule
Entidad & Atributo & Dominio y caracterización \\
\midrule
Juego & \texttt{idJuego} & Identificador único; entero positivo. \\
 & \texttt{nombre} & Cadena de texto no vacía. \\
 & \texttt{costo} & Cantidad monetaria en MXN, positiva y múltiplo de diez. \\
\midrule
Tipo de juego & \texttt{idTipo} & Identificador único; entero positivo. \\
 & \texttt{nombre} & Nombre del tipo de juego; el dominio conserva los tipos
   indicados en el caso de uso. \\
 & \texttt{reglaPuntuacion} & Cadena que describe la regla para asignar puntos
   al jugar ese tipo. \\
 & \texttt{condicion} & Cadena que describe la condición bajo la que se aplica
   la regla de puntuación. \\
\midrule
Máquina & \texttt{idMaquina} & Identificador único; entero positivo. \\
 & \texttt{estado} & Enumeración: \texttt{activo}, \texttt{fuera de servicio}
   o \texttt{mantenimiento}. \\
\midrule
Partida & \texttt{idPartida} & Identificador único; entero positivo. \\
 & \texttt{puntajeObtenido} & Número entero no negativo obtenido por el
   jugador. \\
 & \texttt{puntosGanados} & Número entero no negativo otorgado por la partida,
   de acuerdo con el tipo de juego. \\
 & \texttt{bonificaciónVIP} & Atributo derivado; puntos adicionales que
   corresponden al diez por ciento de la recompensa cuando el cliente tiene la
   categoría VIP vigente. \\
 & \texttt{fecha} & Fecha en la que terminó la partida. \\
 & \texttt{hora} & Hora en la que terminó la partida. \\
\midrule
Mantenimiento & \texttt{idMantenimiento} & Identificador único; entero positivo. \\
 & \texttt{fecha} & Fecha en la que se realizó el mantenimiento. \\
 & \texttt{tipoMantenimiento} & Cadena que identifica el tipo de mantenimiento. \\
 & \texttt{descripciónProblema} & Cadena de texto que describe el problema
   atendido. \\
\bottomrule
\end{tabularx}
\end{center}

\subsection{Entidad Premio}

\begin{center}
\small
\begin{tabularx}{\textwidth}{@{}>{\bfseries\color{unamblue}}p{0.19\textwidth}p{0.28\textwidth}X@{}}
\toprule
Entidad & Atributo & Dominio y caracterización \\
\midrule
Premio & \texttt{idPremio} & Identificador único; entero positivo. \\
 & \texttt{nombre} & Cadena de texto no vacía. \\
 & \texttt{categoría} & Enumeración por puntos: \texttt{bajo} de 20 a 1,000;
   \texttt{medio} de 1,001 a 3,999; y \texttt{grande} de 4,000 a más de
   10,000 puntos. \\
 & \texttt{rangoEdad} & Enumeración de los rangos \texttt{infantil} (3--12),
   \texttt{juvenil} (3--17) y \texttt{adulto} (18 o más). Los rangos se
   conservan tal como aparecen en el caso de uso. \\
 & \texttt{valorAproximado} & Cantidad monetaria aproximada en MXN. \\
 & \texttt{puntosNecesarios} & Número entero de puntos; debe permitir un canje
   a partir de veinte puntos y es consistente con \texttt{categoría}. \\
 & \texttt{cantidadDisponible} & Número entero no negativo que representa la
   cantidad disponible del premio. \\
\bottomrule
\end{tabularx}
\end{center}

La cantidad disponible se conserva como atributo de \texttt{Premio}, tal como
aparece en el diagrama. Su interpretación operativa y su relación con la
sucursal se precisan junto con las relaciones y las decisiones de diseño.

\subsection{Empleado y subentidades operativas}

\begin{center}
\small
\begin{tabularx}{\textwidth}{@{}>{\bfseries\color{unamblue}}p{0.19\textwidth}p{0.28\textwidth}X@{}}
\toprule
Entidad & Atributo & Dominio y caracterización \\
\midrule
Empleado & \texttt{CURP} & Identificador único de la persona empleada. \\
 & \texttt{RFC} & Cadena con formato de RFC. \\
 & \texttt{nombreCompleto} & Atributo compuesto por \texttt{nombre},
   \texttt{apellidoPaterno} y \texttt{apellidoMaterno}; cadenas de texto. \\
 & \texttt{fechaNacimiento} & Fecha válida de nacimiento. \\
 & \texttt{teléfono} & Atributo multivalorado; conjunto de cadenas con formato
   de número telefónico. \\
 & \texttt{correoElectrónico} & Atributo multivalorado; conjunto de cadenas con formato
   de correo electrónico. \\
 & \texttt{fechaContratación} & Fecha válida de contratación. \\
 & \texttt{salario} & Cantidad monetaria no negativa en MXN. \\
\midrule
Gerente & \texttt{cédulaProfesional} & Cadena que identifica la cédula
   profesional del gerente. Hereda \texttt{CURP} y los atributos de
   \texttt{Empleado}. \\
\midrule
Cajero & \texttt{horarioTrabajo} & Enumeración: \texttt{7:00--15:00} o
   \texttt{15:00--23:00}. Hereda \texttt{CURP} y los atributos de
   \texttt{Empleado}. \\
\midrule
Encargado de premios & \texttt{horarioTrabajo} & Enumeración: \texttt{7:00--15:00}
   o \texttt{15:00--23:00}. Hereda \texttt{CURP} y los atributos de
   \texttt{Empleado}. \\
\midrule
Técnico & \texttt{áreaEspecialidad} & Cadena que identifica el área de especialidad
   del técnico. \\
 & \texttt{Certificación} & Atributo multivalorado; conjunto de
   certificaciones. Hereda \texttt{CURP} y los
   atributos de \texttt{Empleado}. \\
\bottomrule
\end{tabularx}
\end{center}

Las subentidades de \texttt{Empleado} no repiten los atributos heredados ni
requieren un identificador distinto: se identifican mediante la \texttt{CURP}
de la superentidad. En la versión actual del diagrama, \texttt{Juego} se
relaciona con \texttt{Tipo de juego}; por ello, los tipos y sus reglas de
puntuación se describen mediante los atributos de esta última entidad y no
como cuatro subentidades independientes.
